\BeleskeID{09_3-s009}
% element: 09_3-s009:e06
% slide-id: 09_3-s009
Cilj je minimum srednje vrednosti kašnjenja pri rastu i opadanju izlaznog napona. Polazno su u svakom invertoru dužine P i N kanala jednake, kao i njihove širine. Širine P kanala zatim menjamo jednako u oba stepena.

U približnom modelu izlaz se opisuje ekvivalentnom otpornošću aktivnog tranzistora i ukupnom kapacitivnošću $C_L$. Za kratke kanale pretpostavlja se zasićenje u delu prelaza koji određuje kašnjenje. Dva prikazana izraza za svaku otpornost predstavljaju alternativne aproksimacije korekcije zbog modulacije dužine kanala.
